\section{Method}
\label{sec:method}

We introduce ReconDrive, a feed-forward framework to process the urban scenes and generate the 4D gaussian splatting in one pass. We start by introducing the problem in Sec.~\ref{sec:method-problem}, followed by our model design in Sec.~\ref{sec:method-gs} and 4D gaussian splatting in Sec.~\ref{sec:method-4dgs}, and finally the training in Sec.~\ref{sec:method-training}.


\subsection{Problem Definition and Overall Inference}
\label{sec:method-problem}
\paragraph{Problem Formulation.} We formalize visual scene reconstruction and novel-view synthesis as the recovery of a continuous 4D representation from discrete observations. The input is an urban scene sequence $D = \{(I_t^o, \mathcal{C}^o, \mathcal{E}_t) \mid t \in [0, T], o=1,\dots,O\}$, where $I_t^o \in \mathbb{R}^{H \times W \times 3}$ is the RGB image captured at time $t$ by the $o$-th camera. Each camera is characterized by a static calibration matrix $\mathcal{C}^o$, which encapsulates both intrinsic parameters and the fixed extrinsic transformation relative to the ego-vehicle coordinate system. The temporal movement of the scene is captured by the time-varying ego-vehicle pose $\mathcal{E}_t \in \mathbb{SE}(3)$, which maps the vehicle's local coordinate frame to a global world frame. The objective is to utilize a feed-forward model to map $D$ to a 4D Gaussian Splatting representation that decouples static geometry from dynamic traffic participants. For any arbitrary timestamp $t' \in [0, T]$ and a novel viewpoint $\mathcal{E}_{t'}^{\text{novel}}$, the model retrieves and projects the corresponding temporal Gaussian kernels to render high-fidelity observations.

\paragraph{Challenges.} Developing a feed-forward 4DGS framework for autonomous driving entails three primary hurdles: 
1) Scalability \textit{vs.} Efficiency: Balancing extensive spatial coverage of large-scale urban environments with the requirements for real-time rendering.
2) Spatio-Temporal Modeling: Reconstructing dynamic scenes where 4D Gaussians must explicitly decouple and model motion, moving beyond the static assumptions of traditional multi-view foundation models.
3) Geometric-Photometric Alignment: Ensuring that regressed Gaussian attributes maintain both accurate 3D geometry for spatial consistency and high photometric fidelity for realistic sensor simulation.
In this work, we address these requirements by extending the capabilities of existing 3D foundation models. We propose a specialized architecture and training strategy that leverages the geometric priors of the foundation model while adapting them to the high-dynamic and calibrated nature of driving scenes.

\paragraph{Inference Pipeline Overview.} To maintain computational tractability over long-duration sequences, we adopt a segment-wise representation. The scene is partitioned into temporal segments $\{[T_{s}, T_{s+1}]\}$, with 4D Gaussians generated for each segment independently. This approach maximizes environmental coverage while limiting the active Gaussian count during rendering to ensure real-time performance. Within each segment, a 4D Gaussian kernel $\mathcal{G}_i(t)$ is defined by a center-moving paradigm:\begin{equation}\mathcal{G}i(t) = \left( \mu_{i}(t), R_{i}, s_{i}, \alpha_{i}, c_{i} \right)\end{equation}where $R_{i} \in \mathbb{SO}(3)$ is the rotation matrix, $s_{i} \in \mathbb{R}^3$ is the anisotropic scaling vector, $\alpha_{i} \in [0, 1]$ is the opacity, and $c_{i} \in \mathbb{R}^{k}$ denotes the spherical harmonic coefficients for view-dependent color. We further employ a static-dynamic composition strategy. For static background Gaussians, the center $\mu_{i}(t)$ remains constant. For dynamic object Gaussians, we assume local linear motion within a segment:\begin{equation}\mu_{i}(t) = \mu_{i, \text{init}} + v_{i} \cdot (t - T_s)\end{equation}where $v_i \in \mathbb{R}^3$ represents the regressed velocity vector. By using the frames $\{T_{s}, T_{s+1}\}$ as context, our ReconDrive framework regresses these spatio-temporal parameters to construct a unified 4D representation. The detailed architecture and training strategy follow in the subsequent sections.



\subsection{Backbone and Prediction Heads}
\label{sec:method-gs}
We employ a feature backbone to extract dense image tokens, which are processed by a dual-path architecture consisting of the Gaussian Center Prediction Head (GCPH) and the Gaussian Parameter Prediction Head (GPPH). Together, these heads regress spatial coordinates and appearance attributes to form the complete Gaussian parameters.

\paragraph{Feature Backbone.}
The input consists of multi-view images from two context frames ($2\times O$ images, where $O=6$ for the nuScenes dataset~\cite{caesar2020nuscenes}). The backbone follows the VGGT~\cite{wang2025vggt} architecture, utilizing two cascaded stages to extract and fuse spatio-temporal features:
\begin{itemize}[leftmargin=9pt]
    \item \textbf{Image Tokenization:} Each image is first downsampled to a resolution of $H' \times W'$ (with a scale factor of 3). We then utilize a DINOv2~\cite{oquab2023dinov2} encoder to transform the resized images into dense feature tokens. The encoder applies a patch size of $p=14$, resulting in tokens of dimension $d=1024$ and a spatial grid resolution of $(H' / p) \times (W' / p)$.
    \item \textbf{Alternating-Attention Fusion:} The tokens from all $2 \times O$ images are concatenated and passed through 24 Transformer layers utilizing Alternating-Attention. This mechanism alternates between frame-wise self-attention, which maintains local structural consistency within each timestamp, and global self-attention, which captures cross-view and temporal geometric correlations.
\end{itemize}

\paragraph{Hybrid Gaussian Prediction Head.} While VGGT~\cite{wang2025vggt} excels at recovering explicit 3D geometry, such as depth and point maps, its latent features are optimized for structural consistency rather than the fine-grained photometric details required for high-fidelity rendering. To bridge this gap, we move beyond direct regression from backbone features and design the Hybrid Gaussian Prediction Head to leverage the geometric priors of 3D foundation models while resolving their inherent limitations. Specifically, it overcomes photometric deficiency by fusing raw image cues into the attribute regression path for high-fidelity rendering. Simultaneously, it eliminates spatial misalignment by explicitly incorporating sensor calibration into the GCPH, ensuring the reconstructed scene is precisely anchored within the ego-vehicle’s calibrated coordinate system. In the Appendix, we further illustrate the spatial misalignment encountered when directly employing VGGT for autonomous driving scene reconstruction and 3D geometry generation.

\begin{itemize}[leftmargin=9pt]
    \item \textbf{Gaussian Center Prediction Head.} We firstly adopt a Dense Prediction Head (DPT)~\cite{ranftl2021vision} to upsample the fused transformer features back to the originally resized image resolution ($H^{'}\times W^{'}$) as $FC_{t}^{o}$. Each $FC_{t}^{o}$ is then input into a $3~\times 3$ convolutional layer to generate the pixel-wise depth map $DP_t^{o}$ with $H^{'}\times W^{'}$ size. Next, we project each depth map into the 3D space with camera intrinsic and extrinsic parameters, to obtain the gaussian 3D center corresponding to the pixel. Note that the two frames' gaussian 3D centers are located in their ego coordinate separately. Utilizing the calibrated sensor data improves more accurate gaussian location prediction.

    \item \textbf{Gaussian Parameter Prediction Head.} Similar to Gaussian Center Prediction Head, we also use DPT to upsample the fused transformer features back to the originally resized image resolution ($H^{'}\times W^{'}$) as $F_{t}^{o}$. Then we incorporate a shortcut connection to concatenate the original image with the upsampled features to $FP_{t}^{o}$. This is extremely important in the gaussian parameters prediction, because it enables the capture of high-frequency texture and color details that may be lost in transformer feature downsampling. A convolutional layer is used to process $FP_{t}^{o}$ to generate the extra gaussian parameters.
\end{itemize}

Finally, we obtain a pixel-wise 3D gaussian set of the two context frames $T_s$ and $T_{s+1}$: $\mathcal{G}(T_s)$ and $\mathcal{G}(T_{s+1})$.

Furthermore, to mitigate the domain gap between general 3D data and large-scale urban driving scenes, we adopt a parameter-efficient fine-tuning strategy. We utilize frozen pre-trained VGGT~\cite{wang2025vggt} weights as a robust structural prior and apply LoRA~\cite{hu2022lora} to adapt the model to the autonomous driving domain. For our implementation, we configure the LoRA adapters with a rank $r=8$ and a scaling factor $\alpha=32$.


\subsection{4D Gaussian Generation}
\label{sec:method-4dgs}
This part further explains how we obtain 4D Gaussians with static-dynamic composition and generated $ \mathcal{G}_{T_{s}}$ and $ \mathcal{G}_{T_{s+1}}$.


\paragraph{Dynamic Object Mask and Motion Estimation.} The autonomous driving environment can be primarily categorized into two components: static background elements (e.g., sky, roads, buildings, and vegetation) and dynamic traffic participants (e.g., vehicles and pedestrians)—the latter being the main contributors to dynamic scene changes. Here, we estimate the motion flow for the Gaussian kernels corresponding to these dynamic traffic participants.




First, we leverage the foundation model SAM2~\cite{ravi2024sam} to extract instance-level regions of traffic participants (with a particular focus on vehicles) from the two context frames. SAM2 is chosen for its robust mask extraction capability and efficient inference speed, which aligns with our real-time inference requirements. 
Next, we calculate the velocities of dynamic objects in the ego-coordinate system of frame $T_{s}$. The nuScenes dataset~\cite{caesar2020nuscenes} provides 3D bounding box annotations for each object. For each dynamic object, we transform its 3D location at frame $T_{s+1}$ into the ego-coordinate system of frame $T_{s}$, then compute the object’s displacement. Within the short time segment $[T_{s}, T_{s+1}]$, we assume the object undergoes rigid-body motion with a constant velocity. Thus, the object’s velocity is estimated as:
\begin{equation}
    v = \textit{displacement}~/~(T_{s+1} - T_{s})
\end{equation}
\begin{wrapfigure}{r}{0.42\textwidth}
\vspace{-8pt}
\centering
\includegraphics[width=0.40\textwidth]{figures/mask-gaussians-mapping.pdf}
\vspace{-10pt}
\caption{\textbf{Mapping Among Images, Masks, Dense Features, and Gaussian Indices.}
(a) Input image. (b) Dynamic object mask.
(c) Dense features. (d) Gaussian kernels.
Each pixel in the dense feature map generates a corresponding Gaussian kernel,
enabling consistent pixel-wise mapping across all these components.}
\vspace{-10pt}
\label{fig:mask-gaussians-mapping}
\end{wrapfigure}
Note that we can also calculate the object displacement by two-frame gaussian centers for dataset without 3d object annotations.
Through this process, we obtain the dynamic object motion flow $V$ with dimensions $H^{'} \times W^{'} \times 3$, where all pixel regions within the object’s mask share the same velocity. As illustrated in Fig.~\ref{fig:mask-gaussians-mapping}, a pixel-wise spatial mapping exists between the image space and the dense feature space associated with the Gaussian kernels. Thus, for the $i$-th Gaussian kernel at frame $T_{s}$, we derive its 3D center at any time $t$ using the following equation:
\begin{equation}
    \mu_{i}(t) = \mu_{i}(T_{s}) + V(i) \cdot (t - T_{s}), \quad t \in [T_{s}, T_{s+1}].
\end{equation}

\paragraph{Temporal Alignment and Fusion.}
Here we explain how we align and fuse the two context-frame Gaussians $\mathcal{G}(T_s)$ and $\mathcal{G}(T_{s+1})$ into final 4D Gaussians for segment $[T_{s}, T_{s+1}]$ by transforming the frame $T_{s}$'s Gaussians into frame $T_s$'s ego coordinate and time $T_s$. Firstly, we spatially transform the frame $T_{s+1}$'s Gaussians into frame $T_{s}$'s ego coordinate system with the ego transformation matrix from frame $T_i$'s ego vehicle pose $E_{T_{i}}$ to frame $T_{i+1}$'s ego vehicle pose $E_{T_{i+1}}$. We mainly consider transforming the gaussian centers, rotation, and color coefficients parameters.  Then, temporally align the $\mathcal{G}(T_{s+1})$ to time $T_s$ by moving the dynamic object gaussian centers with velocity flow $V$. The two steps cannot be swapped because the velocity field is computed in $T_{s}$'s ego frame. Finally, we get transformed Gaussians $\mathcal{G}^{'}(T_{s+1})$, and concatenate the two Gaussians $\mathcal{G}(T_{s})$ and $\mathcal{G}^{'}(T_{s+1})$. These concatenated Gaussians and the velocity flow comprise the final 4D gaussian splatting.


\subsection{Training}
\label{sec:method-training}
\paragraph{Training Loss.}
To enhance visual fidelity, we adopt perceptual loss using a pre-trained VGG-19 network~\cite{simonyan2014very}, which computes the L1 distance between high-level features of rendered images $\hat{I}$ and real images $I$:
$\mathcal{L}_{\text{percep}} = \left\| \phi(\hat{I}) - \phi(I) \right\|_1$,
where $\phi(\cdot)$ denotes the feature extractor of a specific convolutional layer, encouraging perceptually consistent features. L2 loss is used to constrain pixel-wise intensity differences between rendered and real images: $\mathcal{L}_{l2} = \left\| \hat{I} - I \right\|_2$. Combining them, the rendering loss for supervising photometric rendering is formulated as:
\begin{equation}
\mathcal{L}_{\text{render}} = \lambda_{\text{percep}} \cdot \mathcal{L}_{\text{percep}} + \lambda_{l2} \cdot \mathcal{L}_{l2}
\end{equation}

To stabilize Gaussian center prediction training and improve geometric accuracy, we introduce projection loss to enforce consistency when warping a single source frame $t$ to the reference frame $t'$. It combines a weighted L1 loss (measuring pixel-wise intensity differences between warped predicted images $\hat{I}_{\text{warped}}$ and warped ground-truth images $I_{\text{warped, GT}}$) and a weighted SSIM loss (evaluating structural and perceptual consistency via $1 - \text{SSIM}(\hat{I}_{\text{warped}}, I_{\text{warped, GT}})$). Both components are computed under a valid warping mask, with the final loss defined as:
\begin{equation}
\mathcal{L}_{\text{project}} = \mathcal{L}_{\text{masked}} \left( \lambda_{l1} \cdot \mathcal{L}_{l1} + \lambda_{\text{ssim}} \cdot \mathcal{L}_{\text{ssim}} \right)
\end{equation}
where $\mathcal{L}_{\text{masked}}$ is the mask-weighted loss computation, and $\lambda_{l1}, \lambda_{\text{ssim}}$ are the weight coefficients for L1 and SSIM loss.
More projection loss details are provided in the Appendix.

The norm loss acts as a regularization term for Gaussian scale and opacity parameters:
\begin{equation}
\mathcal{L}_{\text{norm}} = \lambda_{\text{scale}} \cdot \mathbb{E}\left[ \|\text{scale}_\text{maps}\|_2 \right] + \lambda_{\text{opacity}} \cdot \mathbb{E}\left[ |\text{opacity}_\text{maps}| \right]
\end{equation}
where $\text{scale}_\text{maps}$ denotes Gaussian scale parameter maps, and the L2-norm $\|\cdot\|_2$ constrains scale magnitude to avoid extreme values; $\text{opacity}_\text{maps}$ represents Gaussian opacity maps, and the L1-norm enforces sparsity by encouraging opacities to converge to 0 or 1. $\lambda_{\text{scale}}$ and $\lambda_{\text{opacity}}$ are corresponding weight coefficients.

Finally, ReconDrive is trained end-to-end with the total loss:
\begin{equation}
\mathcal{L} = \mathcal{L}_{\text{render}} +  \mathcal{L}_{\text{project}} + \mathcal{L}_{\text{norm}}.
\end{equation}~\label{eq:training loss}
